\documentclass{article}
\usepackage[T1]{fontenc}
\usepackage{lmodern}
\usepackage{graphicx}
\usepackage{amsmath,amssymb}
\usepackage[a4paper,margin=2cm]{geometry}
\usepackage[absolute,overlay]{textpos}
\setlength{\TPHorizModule}{1mm}
\setlength{\TPVertModule}{1mm}
\usepackage[indonesian]{babel}
\usepackage{hyperref}
\setlength{\emergencystretch}{2em}
% unit: Halaman 5

\begin{document}

% === Halaman 5 (python/native) ===

• <G, +>  menyatakan sebuah grup dengan operasi penjumlahan.

• <G, > menyatakan sebuah grup dengan operasi perkalian

Contoh-contoh grup:

1. <R, +> : grup dengan himpunan bilangan riil dengan operasi +


e = 0 dan a’ = –a  

2. <R*, > :  grup dengan himpunan bilangan riil tidak nol (yaitu, R* = R – \{ 0\} ) 

dengan operasi kali ()


e = 1 dan a’ = 1/a = a -1

3.   <Z, +> dan <Z, > masing-masing adalah grup dengan himpunan bilangan bulat 

(integer) dengan operasi + dan 

5

Rinaldi Munir/II4021 Kriptografi

\end{document}
